\documentclass[a4paper,12pt]{dinbrief} 
\usepackage{ngerman}
\usepackage[utf8]{inputenc}
%\usepackage{eurosym}
%\usepackage[top=4cm,left=2.9cm,right=3cm,bottom=4cm]{geometry}


\begin{document} 
\setbottomtexttop{2cm}
%\nowindowtics
%\nowindowrules
\subject{\textbf{Widerspruch Bescheid 197726 vom 29.05.2026} }
\signature{Dr.~Fabian Schury}
\place{Aurachtal}
\Datum{\today} 
\backaddress{Dr.~Fabian Schury $|$ Reichenfelser Straße 10 $|$ 91086 Aurachtal}
\address{Dr.~Fabian Schury \\ Reichenfelser Straße 10 \\ 91086 Aurachtal}
 
\begin{letter} 
{   Bezirk Mittelfranken \\
    Postfach 617 \\[0.5cm]
    \textbf{\Large 91511 Ansbach} }
\opening{ Sehr geehrte Damen und Herren,}

hiermit lege ich fristgerecht Widerspruch gegen die Bescheinigung vom
29.05.2026 mit dem Zeichen 197726 ein.

Der Zeitraum der Bewilligung ist sachdienlich, allerdings nicht die bewilligte
Stundenanzahl. Mein Sohn Leopold besucht derzeit in der Kindertageseinrichtung
noch die KiTa-Gruppe, in der die Unterstützung durch eine Integrationshilfe mit
15 Wochenstunden angemessen ist, weil dort ein höherer Personalschlüssel
gegeben ist. 

Ab dem 01.09.2026 wird Leopold allerdings in der selben Einrichtung in die
Kindergartengruppe wechseln. Dort ist eine Unterstützung mit 15 Stunden wie
laut Ihres Bescheids bewilligt leider nicht mehr ausreichend. Ich beantrage
daher in Absprache mit der Einrichtung eine Erhöhung der Stundenzahl auf 28 bis
30 Wochenenstunden. Ein separates Schreiben von der Leitung Fr. Ploner und des
Fachdienstes Fr. Reidy ist Ihnen bereits zugegangen. 

Ich bitte daher um eine entsprechende Neubwertung und Erhöhung der
Wochenstundenanzahl, damit die Betreuung meines Sohnes auch im Kindergarten
gewährleistet ist.

\closing{Mit freundlichen Grüßen} 
\end{letter}
\end{document}
